\section{From Algorithms to Deployment}

% Shared colour palette for this lecture's hand-built diagrams.
% (No theme colour macros exist in the preamble; we use plain xcolor names
%  plus this small section-local palette throughout the new TikZ.)
\definecolor{sysblue}{RGB}{189,215,238}   % model / policy boxes
\definecolor{mlgreen}{RGB}{169,208,142}   % positive / verifiable
\definecolor{boxgreen}{RGB}{197,224,180}
\definecolor{boxblue}{RGB}{180,199,231}
\definecolor{boxred}{RGB}{244,177,177}    % warning / fragile
\definecolor{boxtan}{RGB}{248,203,173}    % human / preference inputs
\definecolor{accentred}{RGB}{192,0,0}     % punchline callouts

% Source/citation footnote with no visible mark (suppresses the stray "0"
% that \footnote[0] leaves behind). Counter is restored so numbering is unaffected.
\newcommand{\srcnote}[1]{%
  \begingroup
  \renewcommand{\thefootnote}{}%
  \footnote{#1}%
  \addtocounter{footnote}{-1}%
  \endgroup}

% ---------------------------------------------------------------- Recap opener
\begin{frame}{Where We Are: Eight Days of Algorithms}
\begin{itemize}\itemsep0.5em
    \item Days 1--8 built the \emph{machinery}: value-based methods, policy
          gradients, actor--critic, continuous control, exploration, model-based
          RL, and offline RL.
    \item \textbf{Day 8 closed on one spine:} \emph{optimism} explores,
          \emph{pessimism} stays safe; the sign of the uncertainty bonus was
          the whole story. Offline RL learned a policy from a \emph{fixed dataset}
          with no environment interaction.
    \item Every method so far assumed one thing for free: \textbf{a reward
          function we could write down.}
\end{itemize}
\vfill
\begin{center}
\itshape Today: where does all of this actually get deployed, and what breaks
when the reward itself is the hard part?
\end{center}
\end{frame}

% ---------------------------------------------------------------- Today's map
\begin{frame}{RL in the Real World: Three Threads}
\begin{itemize}\itemsep0.6em
    \item \textbf{RLHF (the focus).} When reward is hard to specify, \emph{learn}
          it from human (or AI, or verifiable) feedback. This is how modern
          LLMs are aligned, the deep dive for today.
    \item \textbf{Multi-Agent RL (overview).} What changes when other learning
          agents share your environment and the ground keeps shifting underneath
          you.
    \item \textbf{Robotics (overview).} What breaks when the environment is
          physical: sample cost, safety, resets, and the sim-to-real gap.
\end{itemize}
\end{frame}

% ---------------------------------------------------------------- Learning outcomes
\begin{frame}{Learning Outcomes}
\textbf{RLHF: the deep dive.} By the end you should be able to:
\begin{itemize}\itemsep0.2em\small
    \item explain the RLHF pipeline: \textbf{SFT $\rightarrow$ reward model
          $\rightarrow$ PPO/GRPO}, and the role of the KL-to-reference penalty;
    \item turn human comparisons into a trainable loss with \textbf{Bradley--Terry},
          and say why the preference model needs a \emph{sigmoid} rather than a raw
          score gap;
    \item explain \textbf{reward hacking} (over-optimization) and why the KL leash
          exists;
    \item show that ``reward minus $\beta$-KL'' is maximized by $\pi_{\text{ref}}$
          \emph{tilted} by $e^{\,r/\beta}$, and \textbf{derive} the
          \textbf{RM $\leftrightarrow$ DPO} identity from it, so ``your LM is
          secretly a reward model'';
    \item place \textbf{RLHF $\rightarrow$ RLAIF $\rightarrow$ RLVR} on the
          scaling-oversight spectrum.
\end{itemize}
\vspace{0.5em}
\textbf{MARL \& robotics: conceptual overviews.} Know \emph{why each is hard}
and name 2--3 methods for each.
\end{frame}
